\section{1. Großschreibung}
    Wenn ein Verb zum Nomen wird, schreibt man es immer groß:
    \begin{itemize}
        \item schlafen $\rightarrow$ das Schlafen
        \item lernen $\rightarrow$ das Lernen
    \end{itemize}

\section{2. Der bestimmte Artikel „das“}
    Nominalisierte Verben stehen fast immer mit „das“:
    \begin{itemize}
        \item das Essen
        \item das Trinken
        \item das Arbeiten
        \item das Laufen
    \end{itemize}

\section{3. Wie man ein Verb nominalisiert (Regeln)}
    
    \subsubsection*{Regel A: „das“ + Infinitiv}
        Die einfachste und häufigste Form:
        \begin{itemize}
            \item essen $\rightarrow$ das Essen
            \item sprechen $\rightarrow$ das Sprechen
            \item lesen $\rightarrow$ das Lesen
        \end{itemize}
        Diese Form beschreibt die Tätigkeit an sich.
    
    \subsubsection*{Regel B: Mit festen Ableitungssuffixen (-ung, -er, -e, -nis \dots)}
        
        Viele Verben haben eigene feste Nomen, die nicht einfach der Infinitiv sind.
        
        \paragraph{(1) -ung}
            macht aus einem Verb ein abstraktes Nomen:
        \begin{itemize}
            \item erklären $\rightarrow$ die Erklärung
            \item bewegen $\rightarrow$ die Bewegung
            \item entscheiden $\rightarrow$ die Entscheidung
            \end{itemize}
        
        \paragraph{(2) -er / -erin}
            macht Personen, die etwas tun:
        \begin{itemize}
            \item fahren $\rightarrow$ der Fahrer / die Fahrerin
            \item helfen $\rightarrow$ der Helfer / die Helferin
            \end{itemize}
        
        \paragraph{(3) -e}
            häufig bei älteren Wörtern:
        \begin{itemize}
            \item fahren $\rightarrow$ die Fahrt
            \item schließen $\rightarrow$ die Schließung / der Schluss
            \item sagen $\rightarrow$ die Sage
            \end{itemize}
        
        \paragraph{(4) -nis}
            Ergebnis eines Vorgangs:
        \begin{itemize}
            \item erkennen $\rightarrow$ die Erkenntnis
            \item ereignen $\rightarrow$ das Ereignis
            \end{itemize}
        
        \paragraph{(5) Unregelmäßige / eigene Nomen}
            Einige Verben haben komplett eigene Formen:
        \begin{itemize}
            \item gehen $\rightarrow$ der Gang
            \item stehen $\rightarrow$ der Stand
            \item denken $\rightarrow$ der Gedanke
            \end{itemize}

\section{4. Wann benutzt man welche Form?}
    
    \begin{tabular}{|l|l|}
        \hline
        \textbf{Form}   & \textbf{Bedeutung}                    \\
        \hline
         das + Infinitiv & neutrale Beschreibung einer Tätigkeit \\
        -ung            & Ergebnis oder Prozess                 \\
        -er / -erin     & Person, die etwas tut                 \\
        -nis            & Ergebnis eines Vorgangs               \\
        eigene Formen   & historisch gewachsene Nomen           \\
        \hline
    \end{tabular}

\section{Beispiele zum Vergleich}
    \begin{itemize}
        \item das Lernen (die Aktivität)
        \item die Lehre (System / Lerninhalt)
        \item der Lehrer (Person)
        \item das Fahren (Tätigkeit)
        \item die Fahrt (Reise)
        \item der Fahrer (Person)
    \end{itemize}

\subsection{Bildung aus Nomen}
    
    Man bildet Adverbien, indem man eine Präposition mit einem Nomen kombiniert
    (meist im Dativ oder Akkusativ).

\subsection{Beispiele}
    
    \begin{table}[h]
        \centering
        \renewcommand{\arraystretch}{1.3}
        \begin{tabular}{|l|l|p{6cm}|}
            \hline
            \textbf{Beispiel} & \textbf{Bedeutung} & \textbf{Erklärung}            \\
            \hline
            am Anfang         & zu Beginn          & „am“ = „an dem“ + „Anfang“    \\
            \hline
            im Moment         & jetzt              & „im“ = „in dem“ + „Moment“    \\
            \hline
            bei Nacht         & in der Nacht       & traditioneller Ausdruck       \\
            \hline
            mit Absicht       & absichtlich        & wörtlich: „mit einer Absicht“ \\
            \hline
            ohne Zweifel      & zweifellos         & mit „ohne“ + Nomen            \\
            \hline
        \end{tabular}
    \end{table}


\section{Ausdrücke mit \textit{-weise}}

\begin{table}[h]
    \centering
    \renewcommand{\arraystretch}{1.3}
    \begin{tabular}{|l|l|p{6cm}|}
        \hline
        \textbf{Beispiel} & \textbf{Bedeutung}    & \textbf{Herkunft} \\
        \hline
        glücklicherweise  & zum Glück             & aus „Glück“       \\
        \hline
        dummerweise       & leider, auf dumme Art & aus „dumm“        \\
        \hline
        zufälligerweise   & durch Zufall          & aus „Zufall“      \\
        \hline
    \end{tabular}
\end{table}


\section*{Feste adverbiale Nomen}


\section{Feste adverbiale Nomen}

\begin{table}[h]
    \centering
    \renewcommand{\arraystretch}{1.3}
    \begin{tabular}{|l|l|p{6cm}|}
        \hline
        \textbf{Beispiel} & \textbf{Bedeutung} & \textbf{Bemerkung}     \\
        \hline
        abends            & am Abend           & historisch von „Abend“ \\
        \hline
        nachts            & in der Nacht       & --                     \\
        \hline
        tagsüber          & während des Tages  & --                     \\
        \hline
        sonntags          & jeden Sonntag      & --                     \\
        \hline
    \end{tabular}
\end{table}